% RESEARCH_DRAFT_SHA256: f394837515947b566393fd746c2aa77cb7a65ca8705cce140e0c0903fa51124d
% Layout and source-reference wrappers expanded for the web renderer.
\documentclass[11pt]{article}
\usepackage[margin=25mm]{geometry}
\usepackage{fontspec,unicode-math}
\setmainfont{Latin Modern Roman}
\setmathfont{Latin Modern Math}
\usepackage{amsmath,amssymb,mathtools}
\usepackage{xurl,hyperref}
\hypersetup{colorlinks=true,urlcolor=blue,linkcolor=blue}
\setcounter{secnumdepth}{0}
\setlength{\parindent}{0pt}
\setlength{\parskip}{5pt}
\setlength{\emergencystretch}{3em}
\newcommand{\sref}[2]{\hyperlink{source:#1:#2}{[S#2]}}
\newcommand{\cataloguescope}[1]{\paragraph{Source question.}#1}
\newcommand{\floor}[1]{\left\lfloor#1\right\rfloor}
\begin{document}
\section*{Research-attempt provenance}
This standalone research attempt retains the complete original definition body below,
including its historical source and preparation notes. Those notes describe the
earlier catalogue preparation; the new attempt and its verification are identified
separately after that body. No original ID, catalogue status, or source statement is changed.

% UnsolvedMath ID 30005060; source catalogue code OWR-9790365-010
\section{Tight Bounds for Many-User Multiple-Access Coding}\label{30005060}
{\small UnsolvedMath ID 30005060 · Computer Science}
\subsection{Definitions and mathematical statement}

\paragraph{Shared notation.}Standard number systems are denoted by $\mathbb N,\mathbb Z,\mathbb Q,\mathbb R,\mathbb C$; index ranges, computational models and specialized notation follow the cited problem source and the definitions below.


In a Gaussian multiple-access channel, $K_n$ users transmit length-$n$ codewords and the receiver observes their sum plus Gaussian noise. The many-user regime has $K_n/n\to\mu>0$, where $\mu$ is user density. An achievable bound constructs codes meeting the energy-per-information-bit and target-error constraints; a converse rules out performance beyond a bound for every permitted code. In random access, activity of users is itself uncertain and must be accounted for in decoding.
\paragraph{Statement recorded in the supplied source.}
Determine tight achievable and converse bounds for Gaussian many-user multiple-access coding, including the random-access setting. \href{https://doi.org/10.4171/owr/2022/14}{[S2]}
\paragraph{Scope qualification.}
The source compares energy per bit and user density at a fixed target error rate. Payload, power normalization, error metric and activity law remain those of the report; no specific random-access law is added.
\paragraph{Source status.}
The supplied catalogue labels this record open. Its source-status metadata was checked against the supplied snapshot on 2026-10-11. This import does not independently establish current open status or mathematical correctness. The catalogue literature assessment is dated 2026-08-22.
\subsection{Short English statement}
Close the achievable-versus-impossible performance gap for Gaussian channels with linearly many users, including uncertain user activity.
\subsection{Sources}
\begin{itemize}
\item[S1] Ulam AI and the UnsolvedMath Contributors, UnsolvedMath: A Curated Collection of Open Mathematics Problems, record 30005060 (OWR-9790365-010). Catalogue curation and metadata: CC BY 4.0. The supplied source question is retained; definitions and the English summary are editorial additions.
\url{https://huggingface.co/datasets/ulamai/UnsolvedMath}.
\item[S2] New Mathematical Techniques in Information Theory (2023), Oberwolfach Reports, DOI 10.4171/owr/2022/14, pp. 683--707
\url{https://doi.org/10.4171/owr/2022/14}.
\textit{Underlying Oberwolfach Reports material retains its source-specific licence, including CC BY-SA 4.0 where applicable; see THIRD\_PARTY\_NOTICES.md.}
\end{itemize}
\paragraph{Preparation and review.}
Original question attributed to the cited source and UnsolvedMath contributors. Codex (GPT-6; exact serving revision not exposed) assisted with definitions, English summaries, source inspection and formatting. No new human mathematical authorship or independent proof verification is asserted. Canonical ID 30005060 and source code OWR-9790365-010 checked against the complete supplied registry and existing repository statements; known flagged repetitions and obvious identical conjectures excluded. Full mathematical-equivalence/literature review has not been performed. Statement preparation and final review followed contribution prompts C and D; automated checks verify structure and identity, not open status.
\label{end:30005060}

\section{Literature and context check}
The primary report, p.~697, asks to close the inner/outer-bound gap at fixed target error in the many-user regime and then to handle random access. It does not specify a universal activity law in the open-problem paragraph. The calculation below deliberately uses an all-active, independent-message Gaussian model as an auxiliary case. It is a standard converse route rederived here, not a claim to improve the report's bounds.
\section{Research attempt: A per-user Fano converse and the block-error obstruction}
Let $K$ independent messages $W_i$ be uniform on $[M]$, $M\geq2$. User $i$ transmits a deterministic $n$-dimensional vector $X_i(W_i)$ with $\mathbb E\|X_i\|_2^2\leq E$. Let $Y=\sum_iX_i+Z$, with independent $Z\sim N(0,I_n)$. Let $\epsilon_i=\Pr\{\widehat W_i(Y)\ne W_i\}$ and $\bar\epsilon=K^{-1}\sum_i\epsilon_i$. All logarithms in this calculation are natural. Define $h_2(t)=-t\log t-(1-t)\log(1-t)$.

Introduce the error indicator for user $i$. Given $Y$, the true message is determined on the no-error event and has at most $M-1$ possibilities on the error event. Entropy chain rules therefore give
\[
 H(W_i\mid Y)\leq h_2(\epsilon_i)+\epsilon_i\log(M-1).
\]
Here conditioning reduces entropy because the difference is conditional mutual information, an average of nonnegative relative entropies; nonnegativity follows by summing $-\log t\geq1-t$. Summing conditional entropies bounds joint conditional entropy, and concavity of $h_2$ (its second derivative is $-1/[t(1-t)]$) gives
\[
 I(W_1,\ldots,W_K;Y)\geq
 K\{\log M-h_2(\bar\epsilon)-\bar\epsilon\log(M-1)\}.
\]
For an imposed bound $\bar\epsilon\leq\epsilon\leq1-1/M$, the right side with $\epsilon$ in place of $\bar\epsilon$ is still a lower bound, since $h_2(t)+t\log(M-1)$ is increasing up to $1-1/M$.

Independence of the users implies
\[
 \operatorname{tr}\operatorname{Cov}(Y)
 =n+\sum_i\mathbb E\|X_i-\mathbb EX_i\|_2^2\leq n+KE.
\]
To justify the entropy bound directly, compare the density of $Y$ with the Gaussian having the same mean and covariance $s^2I_n$, where $ns^2=\operatorname{tr}\operatorname{Cov}(Y)$. Nonnegativity of relative entropy implies
$h(Y)\leq (n/2)\log(2\pi e s^2)$: integrate the Gaussian log density and use its quadratic expectation $n$. Also $h(Y\mid W_1,\ldots,W_K)=h(Z)=(n/2)\log(2\pi e)$ by translation. Consequently
\[
 K\{\log M-h_2(\bar\epsilon)-\bar\epsilon\log(M-1)\}
 \leq {n\over2}\log(1+KE/n).                 \tag{1}
\]
For $K/n\to\mu$ and fixed $M,E,\epsilon$, (1) gives the necessary condition
\[
 \mu\{\log M-h_2(\epsilon)-\epsilon\log(M-1)\}
 \leq\tfrac12\log(1+\mu E)
\]
in the stated target-error range.

One tempting strengthening applies block-message Fano with the same $\epsilon$. That substitution is invalid. For example, $K$ independent error events of probability $\epsilon\in(0,1)$ have average per-user error $\epsilon$ and block error $1-(1-\epsilon)^K\to1$. This example concerns the relation between error metrics; it is not asserted to arise from every code. A per-user error constraint alone cannot justify replacing the block error by $\epsilon$.
\section{Verification and adversarial audit}
The proof centers the codewords before adding variances; nonzero codeword means cannot create an extra $K^2E$ term. At $E=0$, (1) forbids error below random guessing. The monotonicity restriction on replacing actual average error by a prescribed bound is stated explicitly. Fixed $K$ versus linearly growing $K$ is accounted for before taking limits.
\paragraph{Audit corrections.} Per-user and block error were separated. The auxiliary energy normalization and all-active assumption are declared rather than silently supplied as the original random-access model.
\paragraph{External computation.} None; symbolic arguments and boundary checks only.
\section{Scope and remaining gaps}
The converse is not matched by an achievable construction here, is not asserted sharper than the source bounds, and does not address unknown activity. Closing the many-user performance gap still requires both an adequate converse and achievability under the actual source model.
\paragraph{Integration review.} Codex (GPT-6; exact serving revision not exposed) reviewed the written arguments and expanded layout and source-reference wrappers. Repository metadata and the production web renderer were checked using Python 3.13 and JavaScript with MathJax 3.2.2. This model-assisted review is not independent human proof verification; no PDF compilation or formal proof checking is asserted.

\section{Completion Estimate}
\noindent\textbf{COMPLETION: 0\%}
The calculations below the preserved source statement investigate a concrete route,
but establish no substantiated new progress toward the entire original target.
Known special cases and failed sufficient conditions are not credited as a solution.
\section{References}
The original source references and catalogue attribution are preserved in the
Sources section above. The following URLs were actually consulted for this attempt;
retention of an additional original reference does not claim fresh consultation of it.
\begin{itemize}
\item \url{https://ems.press/content/serial-article-files/46950} (consulted 2026-10-11). 
\end{itemize}
\section{Model and contributor disclosure}
OpenAI gpt-6-astra, requested reasoning configuration ultra, generated this research
attempt on 2026-10-11. The exact serving revision is not exposed. Codex GPT-6
performed orchestration and repository-format review; its exact serving revision
is also not exposed. The mathematical route, boundary checks and final adversarial
self-audit followed contribution prompts A then B in
\texttt{docs/contribution\_prompts.md}, after reading \texttt{CONTRIBUTING.md}.
Python 3.14.2 assembled the standalone documents; any mathematical computation is
specified in Verification. No theorem prover was used. No human mathematical
authorship, substantial human proof correction, or independent human proof checking
is asserted. Model-assisted review, source retrieval, and structural checks do not
establish correctness or novelty. Full private reasoning transcripts are unavailable
as reproducibility artifacts; the complete arguments, inputs and outcomes actually
used are stated here or in the companion batch audit.
\section{Final status}
Unresolved as a new research contribution. This is a concrete route investigation,
not a claim of new complete resolution or a current-status certification. A known
answer reproduced for context does not become a new contribution. The catalogue's
source-status label is preserved; this attempt status does not assert that a
source-labelled solved question is open.
\end{document}
